\section{Alternate Objectives}
\label{sec:alternate-objectives-appendix}

In this section, we briefly discuss the relationship between the small-budget optimal design under our primary impact objective---the improvement in outcomes over two experiment stages---and three alternative objectives given by either (i) assuming only a \emph{single stage}, and/or (ii) maximizing the \emph{probability of passing} the statistical test or tests instead of the improvement in outcomes.
See \Cref{tab:small-budget-alternate-objectives} for the objective definitions.
These objectives serve both to provide a sensitivity analysis of our optimal small-budget design to the specified objective, and to highlight how various objectives of experiment designers or stakeholders may be more or less aligned in terms of their optimal experiment designs.

\begin{table}[h]
  \centering
  \small
  \renewcommand{\arraystretch}{1.2}
  \begin{tabular}{@{} l @{\hspace{0.75em}} c @{\hspace{0.9em}} c @{\hspace{1.0em}} c c c @{}}
    \multicolumn{1}{c}{} &
    \multicolumn{1}{c}{Objective} &
    \multicolumn{1}{c}{Small-budget index} &
    \multicolumn{3}{c}{Components in elliptical decomposition} \\
    \multicolumn{1}{c}{} &
    \multicolumn{1}{c}{$V_Y(\bm{s}_1)$} &
    \multicolumn{1}{c}{$\pi^Y_d$} &
    $\mu_d$ &
    $\cov(\tau_d,\ATE_2)$ &
    $\cov(\tau_d,\ATE_3)$ \\
    \noalign{\vskip 2pt}
    \hline
    \noalign{\vskip 4pt}
    \makecell[l]{Single-stage\\success} &
    $\E\left[\Test_1\right]$ &
    $\dfrac{\E\left[\tau_d\right]}{\sqrt{c_d}}$ &
    $\checkmark$ & & \\[12pt]
    \makecell[l]{Two-stage\\success} &
    $\E\left[\Test_1 \cdot \Test_2\right]$ &
    $\dfrac{\E\left[\tau_d \cdot \Test_2\right]}{\sqrt{c_d}}$ &
    $\checkmark$ & $\checkmark$ & \\[12pt]
    \makecell[l]{Single-stage\\impact} &
    $\E\left[\Test_1 \cdot \ATE_3\right]$ &
    $\dfrac{\E\left[\tau_d \cdot \ATE_3\right]}{\sqrt{c_d}}$ &
    $\checkmark$ & & $\checkmark$ \\[12pt]
    \makecell[l]{Two-stage\\impact} &
    $\E\left[\Test_1 \cdot \Test_2 \cdot \ATE_3\right]$ &
    $\dfrac{\E\left[\tau_d \cdot \Test_2 \cdot \ATE_3\right]}{\sqrt{c_d}}$ &
    $\checkmark$ & $\checkmark$ & $\checkmark$
  \end{tabular}
  \caption{Objective definitions, their small-budget indices, and the mean and covariance-predictivity components that appear when the corresponding coefficient is decomposed under an elliptical prior.
    A checkmark indicates structural dependence in the decomposition, not a fixed sign or a guarantee that the coefficient is nonzero for every parameter value.
  If an aggregate effect is degenerate, the corresponding covariance coordinate is identically zero for all types.}
  \label{tab:small-budget-alternate-objectives}
\end{table}

Analogues of the small-budget optimal pilot design \Cref{cor:optimal-small-budget-pilot} hold under general priors with moment constraints for any of the alternate objectives in \Cref{tab:small-budget-alternate-objectives}: that is, whenever the corresponding small-budget index has a unique strictly positive maximizer, the optimal small-budget design invests the entire budget in a single type given by the index.
The general index formula can be derived from \linkedCref{thm:optimal-small-budget-design-general-appendix}.
Further, when the prior is elliptical with finite second moments, these small-budget indices can be represented as cost-adjusted combinations of the type prior mean and the type's predictivity for the follow-up and target-population average effects.
\Cref{cor:named-objective-elliptical-decompositions-appendix} gives the coefficient decompositions underlying the checkmarks in \Cref{tab:small-budget-alternate-objectives}.
The checkmarks give the structure of the result---at special parameter values a listed component can have zero coefficient.

Note that the two-stage impact objective we focus on is the only one whose elliptical-prior decomposition involves both follow-up-effect predictivity and target-population-effect predictivity.
These two covariance terms are equal when the follow-up is representative, so in this case two-stage impact is similar to both the single-stage impact and two-stage success objectives, although the weights are different.
Meanwhile, optimizing for passing the pilot statistical test only depends on each type's cost and prior mean
